% =====================================================================
% Springer LNCS (llncs.cls) 会议论文集最小可编译骨架
% 编译: latexmk -pdf springer_llncs.tex          (pdfLaTeX)
% 需要 llncs.cls (CTAN: llncs) 与 splncs04.bst
% 关键: \titlerunning / \authorrunning 作者多时必填, 否则页眉溢出
% 参考文献必须用 splncs04, 不要自选样式
% =====================================================================
\documentclass[runningheads]{llncs}

\usepackage{graphicx}
\usepackage{amsmath}
\usepackage[T1]{fontenc}
\usepackage{hyperref}
\renewcommand\UrlFont{\color{blue}\rmfamily}

\begin{document}

\title{Paper Title Here\thanks{Optional funding/acknowledgement}}
\titlerunning{Short Title}        % 页眉短标题

\author{First Author\inst{1}\orcidID{0000-0000-0000-0000} \and
Second Author\inst{2}}
\authorrunning{F. Author et al.}  % 页眉短作者

\institute{First Institution, City, Country \\
\email{first@example.com} \and
Second Institution, City, Country \\
\email{second@example.com}}

\maketitle

\begin{abstract}
This is the abstract. State problem, method, results and contributions in one paragraph.
\keywords{Keyword1 \and Keyword2 \and Keyword3}
\end{abstract}

\section{Introduction}
Body text. Citation example~\cite{ref1}. See Fig.~\ref{fig:arch} and Eq.~(\ref{eq:main}).

\section{Method}
\begin{equation}\label{eq:main}
  y = f(x) + \epsilon
\end{equation}

\begin{figure}[t]
  \centering
  % \includegraphics[width=0.6\textwidth]{fig1.pdf}
  \rule{0.6\textwidth}{3cm}   % 占位框, 有图时替换
  \caption{System architecture.}
  \label{fig:arch}
\end{figure}

\section{Conclusion}
Summary.

% --- 参考文献: Springer 要求 splncs04 -----------------------------
\bibliographystyle{splncs04}
% \bibliography{refs}
% 演示用内联:
\begin{thebibliography}{1}
\bibitem{ref1}
Author, A.: Title of paper. In: Proc. of Conf., pp. 1--10 (2024)
\end{thebibliography}

\end{document}
